\begin{corollary}[5.10.9]
\label{IV.2.s5.10.9-fr}
Si un anillo local noetheriano $A$ verifica $(S_2)$ y es catenario, es
equidimensional.
\end{corollary}

La hipótesis de (5.10.8) se verifica entonces para
$X=\operatorname{Spec}(A)$, con $d=2$. Para demostrar que todas las
componentes irreducibles de $X$ tienen la misma dimensión, basta entonces, en
virtud de (5.10.8), demostrar que dos de tales componentes $X'$, $X''$ tienen
la misma dimensión cuando se supone además que
$\operatorname{codim}(X'\cap X'',X)=1$. Existe entonces una componente
irreducible $Z$ de $X'\cap X''$ tal que
$\operatorname{codim}(Z,X)=1$, por tanto $\operatorname{codim}(Z,X')=1$, pues
$\operatorname{codim}(Z,X)\geq\operatorname{codim}(Z,X')\geq1$; análogamente,
$\operatorname{codim}(Z,X'')=1$, y como $X$ es catenario, esto implica
$\dim(X')=\dim(X'')$.

\begin{proposition}[5.10.10]
\label{IV.2.s5.10.10-fr}
Sean $X$ un preesquema localmente noetheriano, $Z$ una parte de $X$ estable
por especialización, $\mathcal F$ un $\mathcal O_X$-módulo coherente y
supongamos que el $\mathcal O_X$-módulo
$\mathcal H^0_{X/Z}(\mathcal F)$ es coherente. Entonces:
\begin{enumerate}
\item[\textup{(i)}] Se tiene
$\operatorname{prof}_Z(\mathcal H^0_{X/Z}(\mathcal F))\geq2$.
\item[\textup{(ii)}] Para todo punto
$x\in\operatorname{Ass}(\mathcal F)\cap(X\setmin Z)$, se tiene
$\operatorname{codim}(Z\cap\overline{\{x\}},\overline{\{x\}})\geq2$.
\item[\textup{(iii)}] El conjunto $U$ de los $x\in X$ tales que
$\operatorname{prof}_{Z_x}(\widetilde{\mathcal F}_x)\geq2$ \textup{(con las notaciones
de (5.9.6))} es abierto en $X$; se tiene $X\setmin U\subset Z$, y $U$ es el mayor
abierto de $X$ tal que $\mathcal F|U$ sea $(Z\cap U)$-cerrado.
\end{enumerate}
\end{proposition}

\oldpage[IV-2]{118}

Pongamos, para abreviar, $\mathcal F'=\mathcal H^0_{X/Z}(\mathcal F)$. Se sabe
por (5.9.11) que $\mathcal F'$ es $Z$-cerrado, y la afirmación (i) resulta por
tanto de (5.10.5), aplicado a $\mathcal F'$. Sea
$x\in\operatorname{Ass}(\mathcal F)\cap(X\setmin Z)$; como las restricciones de
$\mathcal F$ y $\mathcal F'$ a $X\setmin Z$ son canónicamente isomorfas (5.9.11),
se tiene $x\in\operatorname{Ass}(\mathcal F')$. Consideremos un punto
$y\in Z\cap\overline{\{x\}}$; el ideal primo $\mathfrak p$ de
$\mathcal O_y$ correspondiente a $x$ está asociado al $\mathcal O_y$-módulo
$\mathcal F'_y$; por tanto, en virtud de (i) y de (\textbf 0, 16.4.6.2), se tiene
$2\leq\operatorname{prof}(\mathcal F'_y)
\leq\dim(\mathcal O_y/\mathfrak p)
=\operatorname{codim}(\overline{\{y\}},\overline{\{x\}})$,
de donde se sigue (ii). Por último, para probar (iii), observemos que $U$ es
el conjunto de los $x\in X$ tales que $\widetilde{\mathcal F}_x$ sea
$Z_x$-cerrado (5.10.5), o también, en virtud de (5.9.6), el conjunto de los
puntos en los que el homomorfismo canónico
$\mathcal F_x\to\mathcal F'_x$ es biyectivo; es, por tanto, el complemento de
la unión de los soportes de $\operatorname{Ker}(\rho_{X/Z})$ y
$\operatorname{Coker}(\rho_{X/Z})$, y ambos son
$\mathcal O_X$-módulos coherentes en virtud de la hipótesis (\textbf 0$_\mathrm I$,
5.3.4) y tienen, por consiguiente, soporte cerrado (\textbf 0$_\mathrm I$, 5.2.2);
esto demuestra que $U$ es abierto, y $U$ es evidentemente el mayor abierto tal
que $\mathcal F|U$ sea $(Z\cap U)$-cerrado; finalmente, la inclusión
$X\setmin U\subset Z$ resulta de (5.9.11).

Veremos más adelante (5.11.1) que, en los casos más importantes, la
afirmación (ii) implica recíprocamente que
$\mathcal H^0_{X/Z}(\mathcal F)$ es coherente.
